\answer{ex:05:1}{$\ell\approx17\um$; $\approx100$~ngày: việc phát đi khắp não là nhờ dây dẫn phân nhánh, khuếch tán chỉ làm nhòe mỗi vị trí phóng ra khoảng $\sim20\um$.}
\answer{ex:05:2}{$r=ab/(1+G)$, $S=1/(1+G)$ chính xác; $+1.7\,\%$ thay vì $+20\,\%$.}
\answer{ex:05:3}{$T^*=1.21\,\text{s}$ khi $G=2$ và $0.19\,\text{s}$ khi $G=9$: thực tế $G\sim2$--$4$, sai lệch được nén $3$--$5$ lần.}
\answer{ex:05:4}{$\mathrm{CV}=1/\sqrt{2\lambda\tau_c}\approx9\cdot10^{-4}$ với tổng tần số $\lambda$; một nơ-ron cần $\tau_c=12.5\,\text{s}$.}
\answer{ex:05:5}{$c\approx0.40$ trong cả hai trường hợp; $\mathrm{CV}=0.050$ so với $0.158$, giảm $\sqrt{10}$ lần. Tần số của từng nơ-ron vẫn tỷ lệ với $a_i$: chỉ có tổng được điều chỉnh.}
\answer{ex:05:6}{$N_1=\overline{A\vee B}$, $N_2=\overline{A\vee N_1}$, $N_3=\overline{B\vee N_1}$, $N_4=\overline{N_2\vee N_3}$, XOR $=N_5=\overline{N_4}$; mỗi lần chuyển trạng thái mất $\tau_c\ln2$, đường qua $4$ cổng --- $2.8\,\text{s}$ cho một XOR.}
\answer{ex:05:7}{(a)~$\tau_\gamma=6/\ln3\approx5.5\,\text{s}$. (b)~$\bar D<4\,\text{s}$ so với $D<2.2\,\text{s}$: độ trễ khó đoán làm con vật kiên nhẫn hơn.}
\answer{ex:05:8}{$k_2=1/\tau_d=10\,\text{s}^{-1}$, $k_1=1/\tau_d-1/\tau_\gamma=9.5\,\text{s}^{-1}$; $\tau_\gamma=1/(k_2-mk_1)$: tăng gấp đôi khi $m$ tăng $+2.6\,\%$ (còn khi $+5.3\,\%$ thì tầm nhìn vô hạn).}
